\section{Sampling at a prescribed irrational angle}\label{sec:parameters}

Fix $X=\R/\Z$ with Haar probability $\mu$, and put
$\sigma x=x+\theta$, where $\theta\in(0,1)$ is a fixed irrational number.
We write $\{t\}\in[0,1)$ for fractional part.  Our immediate aim is to
construct rare logarithmic drops whose characters behave almost like
independent coordinates in finitely many sampling tests.  The frequencies
will be divisible integers; the angle $\theta$ will never be changed.

We use the following elementary facts about irrational rotations.
For each nonzero integer $a$, the average of $e^{2\pi ian\theta}$ over
$1\le n\le N$ tends to zero, by the formula for a geometric sum.
Trigonometric approximation gives equidistribution against continuous
functions, and continuous upper and lower bounds give it for intervals.
The same argument gives convergence uniform in the starting point for
continuous functions.  In particular, every positive orbit is dense.
These facts also hold with $\theta$ replaced by any positive integer
multiple of $\theta$.

For an integer $m\ge200$, put $h=m^{-3/5}$ and
$\Delta(z)=\min(\{z\},1-\{z\})$.  Define on two circles
\begin{equation}\label{eq:sampling-kernel}
 w_m(z,v)=\mathbf1_{\{\Delta(z)\ge m^{-2}\}}
       \frac{\mathbf1_{\{0<v<h\Delta(z)\}}}{h\Delta(z)},
 \qquad
 w_m^{(1)}(z,v)=w_m(z,v),\quad
 w_m^{(2)}(z,v)=w_m(z-v,v),
\end{equation}
where $v$ is represented in $[0,1)$ and the expression is zero when the
cutoff fails.  The first kernel forces a small positive move with no
wrap around zero.  The second is its density viewed from the endpoint.

The kernels have almost unit mass at either endpoint.  We next control
their second moments and the correlations produced by two separated
integer powers.  These estimates will let us average products over a
whole block of coordinates.  The scalar lemmas also serve the
independent-coordinate construction in Appendix~\ref{p:sec:product}.

\begin{lemma}[Scalar sampling estimates]\label{lem:scalar-sampling-estimates}
Let $m\ge200$ be an integer, put $h=m^{-3/5}$, and let
$u=w_m^{(r)}$ for either $r=1$ or $r=2$.  Write
$\bar u(s)=\int_0^1u(s,v)\,dv$.  Uniformly in $s$, the supremum plus
total variation of $v\mapsto u(s,v)$ is $O(m^{13/5})$, hence
$O(m^{10})$, and
\begin{equation}\label{eq:kernel-estimates}
 \begin{aligned}
 &\|\bar u\|_\infty\le2,\qquad
   \int_0^1\bar u=1-2m^{-2},\\
 &(1-2m^{-2})^2\le\int_0^1\bar u^2
       \le1+\frac{h^2}{(1-h)^2},\qquad
   \iint_{[0,1]^2}u^2=\frac2h\log\frac{m^2}{2}.
 \end{aligned}
\end{equation}

For the sampling estimate let $0<N_0\le N\le m^2$,
$Q=m^{(7/10)m}$ and $H=m^{300}$.  Let $\lambda$ be uniform on
$[1/N_0,2/N_0]$.  If $n\in[Q/2,Q]\cap\N$, the density $p_n$ of
$\{n\lambda\}$ satisfies
\begin{equation}\label{eq:scalar-marginal}
 |p_n-1|\le N_0/n\quad\hbox{almost everywhere}.
\end{equation}
If two distinct such powers $n,n'$ satisfy
\[
 |vn+v'n'|>Nm^{1000}
 \quad(0<|v|+|v'|,\ |v|,|v'|\le H,\ v,v'\in\Z),
\]
then, for intervals $I,J\subset[0,1)$,
\begin{equation}\label{eq:scalar-rectangles}
 \left|\mathbb P(\{n\lambda\}\in I,\{n'\lambda\}\in J)
             -|I||J|\right|=O(m^{-80}).
\end{equation}
Consequently, for complex functions $f,g$ of bounded variation on
$[0,1]$, with $B(f)=\|f\|_\infty+\operatorname{Var}(f)$,
\begin{equation}\label{eq:scalar-bv}
 \left|\mathbb E[f(\{n\lambda\})g(\{n'\lambda\})]
           -\left(\int_0^1f\right)\left(\int_0^1g\right)\right|
       \le C m^{-80}B(f)B(g).
\end{equation}
All implicit constants are absolute, for sufficiently large $m$,
and uniform in $N_0,N,n,n'$ and the intervals or functions.
\end{lemma}

\begin{proof}
For the first orientation the function of $v$ is a constant height
at most $m^{13/5}$ on an interval, so the variation bound is immediate.
For the second, nonzero values require $v<s$ in representatives
$s,v\in[0,1)$: if $v>s$, then $\Delta(s-v)\le v-s<v$, whereas
$h<1$.  On $0<v<s$, the conditions
$\Delta(s-v)\ge m^{-2}$ and $v<h\Delta(s-v)$ describe an interval.
The reciprocal of $\Delta(s-v)$ has at most two monotone pieces
there, with the same height bound.  This proves the asserted uniform
bound in both orientations.

In the first orientation $\bar u$ is the cutoff indicator, and
\[
 \iint u^2=\frac2h\int_{m^{-2}}^{1/2}\frac{ds}{s}
          =\frac2h\log(m^2/2).
\]
The change $(s,v)\mapsto(s-v,v)$ preserves double integrals, giving
the same total mass and square integral for the second orientation.
To bound its averaged density, first omit the cutoff.  The resulting
density $d$ is the law of $z+t\Delta(z)$ for independent uniform
$z\in[0,1)$ and $t\in(0,h)$.  For fixed $t$ this is an increasing
piecewise-linear bijection of $[0,1]$, with slopes $1+t$ and $1-t$.
Thus $\int d=1$ and $|d-1|\le h/(1-h)$, whence
\[
 \int d^2=1+\int(d-1)^2\le1+\frac{h^2}{(1-h)^2}.
\]
Restoring the cutoff only decreases this nonnegative density;
also $d\le(1-h)^{-1}\le2$.  These observations give the upper
bounds in \eqref{eq:kernel-estimates}, while Cauchy--Schwarz gives
the lower bound from its exact mass.  The cancellation of the linear
term here is essential when multiplying estimates over $m$ coordinates.

Periodizing a uniform interval of length $n/N_0$ gives
\eqref{eq:scalar-marginal}; in particular $p_n\le2$ for large $m$.
For the separated powers and any permitted nonzero frequency pair,
direct integration on $[1/N_0,2/N_0]$ gives
\[
 \left|\mathbb E e^{2\pi i(vn+v'n')\lambda}\right|
 \le\frac{N_0}{\pi|vn+v'n'|}\le\frac1\pi m^{-1000}.
\]
To pass from characters to rectangles, convolve each interval
indicator with the Fej\'er kernel
\[
 F_H(s)=\frac1H\left(\frac{\sin(\pi Hs)}{\sin(\pi s)}\right)^2.
\]
Expanding the squared geometric sum shows that $F_H$ is nonnegative,
has mass one, and has frequencies of magnitude below $H$.
The convolved indicators lie in $[0,1]$, preserve their means, and
have Fourier coefficients of modulus at most one.  Outside strips
of width $\eta=m^{-80}$ about interval endpoints, their error is
$O((H\eta)^{-1})$, by integrating the kernel tail.  The marginal
bound $p_n,p_{n'}\le2$ shows that replacing both indicators in a
product costs $O(\eta+(H\eta)^{-1})$, both for the actual joint
law and for Haar product measure.  The remaining Fourier error is
$O(H^2m^{-1000})$.  This proves \eqref{eq:scalar-rectangles}; no
joint density, or independence of the sampled powers, is used.

Finally, a bounded-variation function is, up to a null set, a constant
plus an integral of indicators of intervals $[a,1)$ against its
complex Stieltjes measure.  The constant and measure have total size
at most $B(f)$; the same holds for $g$.  Apply
\eqref{eq:scalar-rectangles} to the resulting products of interval
indicators, allowing the full interval.  Integration against total
variation gives \eqref{eq:scalar-bv}.  Changes at the countably many
possible discontinuities do not matter, because the marginals have
densities.
\end{proof}

Before choosing the block scales, we record the finite-coefficient bound
that will determine how much attenuation each block must supply.

\begin{lemma}[A nonzero finite convolution]\label{lem:polynomial-coefficient}
For every integer $K\ge0$ and $D\ge\delta>0$ there is
$\varepsilon=\varepsilon(K,D,\delta)>0$ such that complex Laurent polynomials
$A,B$ supported in $[-K,K]$, with coefficients of modulus at most $D$
and $|A_0|,|B_0|\ge\delta$, satisfy
\[
 \max_{0\le n\le2K}|(AB)_{-n}|\ge\varepsilon.
\]
\end{lemma}

\begin{proof}
The admissible polynomial pairs form a nonempty compact set.
For each pair, the first nonzero exponents of $A$ and $B$ are at most
zero, and their sum belongs to $[-2K,0]$.  The corresponding product
coefficient is nonzero.  The continuous displayed maximum therefore
has a positive minimum on the compact set.
\end{proof}

We now choose the parameters for the fixed circle rotation.  First prepare
data for the later continuity argument.  Let $O_j$ be the
open arc of radius $1/(10j)$ about zero.  Choose a continuous nonnegative
bump supported in $O_j$ and positive on a smaller arc.  Finitely many of
its translates $x\mapsto\psi(\sigma^n x)$, with $0\le n\le h_j^*$,
have a strictly positive sum on $X$, by density and compactness.
Normalize by this sum to obtain continuous functions $\phi_{j,n}\ge0$
such that
\begin{equation}\label{eq:partition}
 \sum_{n=0}^{h_j^*}\phi_{j,n}=1,
 \qquad \sigma^n(\operatorname{supp}\phi_{j,n})\subset O_j.
\end{equation}
All $h_j^*$ and $\phi_{j,n}$ are fixed before the construction below.

The parameters below are grouped into consecutive blocks of $m_j$
coordinates.  In block $j$, $K_j$ is the cutoff for a finite convolution,
$C_j$ is the cost of each logarithmic drop, and $L_j$ is the attenuation
needed to make the omitted tail smaller than a finite polynomial
coefficient.  The total phase increment in the block is of order
$1/K_j$.  The sampling horizon $Q_j$ lies well below the next cutoff
$K_{j+1}$, so sampling moves can be controlled in every future block.
The two widths $a_j,b_j$ will provide crossings in complementary ranges
of the tail index.  The proposition states these requirements together;
its proof shows that they can all be chosen successively.

\begin{proposition}\label{prop:parameters}
For the fixed $\theta$ there are positive integers $m_j,K_j,G_j,q_i$,
positive real numbers $L_j,C_j,\varepsilon_j$, and consecutive blocks
$J_j=\{d_{j-1}+1,\ldots,d_j\}$ of size $m_j$, with $d_0=0$, having
the following properties.  Set
\begin{gather*}
 q_0=1,\qquad u_i(x)=\{q_ix\},\qquad\alpha_i=\{q_i\theta\},\\
 \kappa_j=(100j)^{-1},\qquad N_j=m_jK_j,\\
 a_j=G_j/m_j,\qquad b_j=m_j^{-19/20},\qquad Q_j=m_j^{(7/10)m_j}.
\end{gather*}
The frequencies are strictly increasing and divisible, and
\begin{equation}\label{eq:phases}
 \frac1{N_j}<\alpha_i<\frac2{N_j},\qquad
 q_i\mid q_{i+1},\qquad \frac{q_{i+1}}{q_i}\ge N_j
 \quad(i\in J_j).
\end{equation}
For every integer $k$ with $N_j/4\le |k|\le m_j^{(4/5)m_j}$, at least
$m_j/20$ indices in $J_j$ satisfy
\begin{equation}\label{eq:spread}
 1-\{k\alpha_i\}\in[b_j,1-b_j].
\end{equation}
If $G_jK_j\le |k|\le N_j/4$, every index in $J_j$ has
$1-\{k\alpha_i\}\in[a_j,1-a_j]$.

The initial values are $K_1=100,G_1=1$.  The numbers $L_j\ge1$ are
nondecreasing, $C_j=\max(L_j,K_j^2)$, and the sizes satisfy
\begin{align}\label{eq:scales}
 &m_j\ge\max(K_j,2^j,4G_j,200),\qquad
 K_{j+1}=\lceil m_j^{(3/4)m_j}\rceil>2K_j,\notag\\
 &G_{j+1}K_{j+1}<m_j^{(4/5)m_j},\qquad
 \kappa_jm_j^{1/40}\ge L_j,\qquad
 \kappa_{j+1}\sqrt{G_{j+1}}\ge L_j,\notag\\
 &h_{j+1}^*/K_{j+1}\le2^{-j}.
\end{align}
Put $D_j=\exp((K_j+1)\sum_{s<j}C_sm_s)$.  If $A,B$ are Laurent
polynomials supported in $[-K_j,K_j]$, every coefficient has modulus at
most $D_j$, and $|A_0|,|B_0|\ge1/j$, then
\begin{equation}\label{eq:polynomial}
 \max_{0\le n\le2K_j}|(AB)_{-n}|\ge\varepsilon_j,
 \qquad
 (2K_j+1)3e^{-L_j/10000}\le\varepsilon_j/j.
\end{equation}

Finally put $p_*=d_{j-1}$ and
$\ell_j=((j+1)^2(1+p_*)q_{p_*})^{-1}$.
For every $t=r/q_{p_*}$ with $r\in\Z$, $0\le r<q_{p_*}$, there is a finite nonempty set
$\mathcal N_{j,t}$ of positive integers in $[Q_j/2,Q_j]$, with
$\{n\theta-t\}\in(0,\ell_j)$, such that for $r'=1,2$,
\begin{equation}\label{eq:sampling}
 \left\|\frac1{|\mathcal N_{j,t}|}\sum_{n\in\mathcal N_{j,t}}
       \prod_{i\in J_j}w_{m_j}^{(r')}(u_i(x),\{n\alpha_i\})-1
       \right\|_{L^2(\mu)}<\frac1j.
\end{equation}
\end{proposition}

The proof separates two choices.  We first choose independent trial
phases to obtain the finite probabilistic tests.  Only afterwards do
we realize them, to sufficient accuracy, by characters of the prescribed
rotation.  Large gaps between successive frequencies also make those
characters converge in distribution to independent Haar variables.

\begin{proof}
Suppose $K_j,G_j$ and all earlier blocks are fixed.  Apply
Lemma~\ref{lem:polynomial-coefficient} with $K=K_j$, $D=D_j$ and
$\delta=1/j$ to choose $\varepsilon_j$.
Choose $L_j\ge\max(1,L_{j-1})$ to satisfy the second inequality in \eqref{eq:polynomial}
(omit $L_{j-1}$ when $j=1$), set $C_j=\max(L_j,K_j^2)$, and then choose
an integer $G_{j+1}$ with $\kappa_{j+1}\sqrt{G_{j+1}}\ge L_j$.
All requirements of \eqref{eq:scales} now hold for sufficiently large
$m=m_j$.  In what follows earlier data remain fixed while $m\to\infty$;
in particular $p_*,q_{p_*},\ell_j$ and the finite list of shifts $t$ are
fixed.  Write $N=mK_j$, $Q=m^{(7/10)m}$, $a=G_j/m$, and
$b=m^{-19/20}$.

\paragraph{Separated sampling powers.}
Equidistribution gives at least a fixed positive multiple of $\ell_j Q$
integers $n\in[Q/2,Q]$ with $\{n\theta-t\}\in(0,\ell_j)$ for
each required $t$, once $m$ is sufficiently large.  Set
$H=m^{300}$ and $R=Nm^{1000}$.
We can select at least $m^{0.69m}$ such powers, separately for every $t$,
so that distinct selected powers obey
\begin{equation}\label{eq:power-separation}
 |vn+v'n'|>R\quad
 (|v|,|v'|\le H,\ (v,v')\ne(0,0),\ v,v'\in\Z).
\end{equation}
Indeed, if $v'\ne0$, each previously chosen $n$ excludes at most
$2R+2$ integers $n'$ for each frequency pair $(v,v')$.
There are at most $(2H+1)^2$ pairs, and $N\le m^2$.
A greedy selection thus loses only a polynomial factor in $m$ from the
available number.  Since $Q=m^{0.70m}$, it eventually leaves at least
$m^{0.69m}$ powers.  If $v'=0$ and $v\ne0$, then $|vn|\ge Q/2>R$
for large $m$, so this case causes no exclusion.  Fix the resulting
sets $\mathcal N_{j,t}$.

\paragraph{Trial phases and the spread condition.}
For $i\in J_j$, take independent uniform $\lambda_i\in[1/N,2/N]$.
For each integer $k$ in the range of \eqref{eq:spread}, the law of
$\{k\lambda_i\}$ has density at most $5$: it is the periodization of
a uniform interval of length $|k|/N\ge1/4$, whose density is at most
$1+N/|k|$.  The probability that an index fails the strict interior
version of \eqref{eq:spread} is therefore at most $10b$.
If fewer than $m/20$ indices succeed, more than $19m/20$ fail.
Independence and a union bound give failure probability at most
$2^m(10b)^{(19/20)m}$.  A further union over both signs of $k$, up to
$m^{(4/5)m}$, tends to zero: the coefficient of $m\log m$ in its
logarithm is $4/5-(19/20)^2<0$.
The lower range in the proposition requires no probabilistic choice.
When $G_jK_j\le|k|\le N/4$, every trial phase has
$a\le|k|\lambda_i\le1/2$; the same observation will hold for the
eventual $\alpha_i$.

\paragraph{One-coordinate and two-power estimates.}
Fix an orientation, put $u=w_m^{(r')}$, and set
$\bar u(s)=\int_0^1u(s,v)\,dv$ and $A_*=\int_0^1\bar u(s)^2\,ds$.
Lemma~\ref{lem:scalar-sampling-estimates}, with $N_0=N$, gives
$\|\bar u\|_\infty\le2$,
$(1-2m^{-2})^2\le A_*\le1+4h^2$ and
$\iint u^2\le4h^{-1}\log m$, where $h=m^{-3/5}$.
Its bounded-variation estimate also gives, uniformly in $s$,
\begin{equation}\label{eq:two-power}
 \mathbb E[u(s,\{n\lambda_i\})u(s,\{n'\lambda_i\})]
       =\bar u(s)^2+O(m^{-50})\quad(n\ne n').
\end{equation}
Indeed, the error from \eqref{eq:scalar-bv} is $O(m^{-60})$.
The one-power density estimate \eqref{eq:scalar-marginal} likewise gives
$\mathbb E[u(s,\{n\lambda_i\})\bar u(s)]
=\bar u(s)^2+O(m^{-50})$, since $n\ge Q/2$.

\paragraph{The product sampling test.}
We now have one-coordinate and two-power estimates.  Their role is to
control the second moment of an average of products over a whole block;
the diagonal terms will be absorbed by the large number of sampling
powers, while the distinct-power terms approach the ideal product law.
For independent Haar variables $z=(z_i)_{i\in J_j}$, and a fixed power
set of size $P$, put
\[
 Z(z)=\frac1P\sum_{n\in\mathcal N_{j,t}}
                 \prod_{i\in J_j}u(z_i,\{n\lambda_i\}),
 \qquad Z^{\rm id}(z)=\prod_{i\in J_j}\bar u(z_i).
\]
By independence of the trial phases over $i$, the distinct-power terms
of $\mathbb E\|Z\|_2^2$, and the terms of
$\mathbb E\langle Z,Z^{\rm id}\rangle$, are products of $m$ factors
$A_*+O(m^{-50})$.  They equal $A_*^m+O(m^{-49})$ by telescoping,
because $(1+4h^2+O(m^{-50}))^m$ stays bounded.
The equal-power terms are each at most $(8h^{-1}\log m)^m$ by
\eqref{eq:kernel-estimates} and the one-phase density bound, and their
total averaging weight is $1/P$.  Thus
\[
 \mathbb E\|Z-Z^{\rm id}\|_2^2
 \le O(m^{-49})+P^{-1}(8h^{-1}\log m)^m\longrightarrow0,
\]
since $P\ge m^{0.69m}$ and $h^{-1}=m^{0.60}$.
Also
\[
 \|Z^{\rm id}-1\|_2^2
 =A_*^m-2(1-2m^{-2})^m+1\longrightarrow0
\]
by \eqref{eq:kernel-estimates} and $mh^2=m^{-1/5}\to0$.
It follows that $\mathbb E\|Z-1\|_2^2\to0$.
There are only finitely many $t$ at this stage.  Markov's inequality,
the two orientations, and the spread estimate therefore let us fix
trial phases with strict spread inequalities and sampling norms below
$1/(2j)$, simultaneously for every test.  We can also avoid the finitely
many equalities $n\lambda_i\in\Z$ and the endpoints $1/N,2/N$, since
each has probability zero.

\paragraph{Realizing the phases by circle characters.}
Now fix this $m$ and these trial phases.  Successively choose $q_i$ as
multiples of the preceding frequency so that $\{q_i\theta\}$ is as
close to $\lambda_i$ as desired and $q_i/q_{i-1}$ is as large as
desired.  This is possible because the positive multiples of
$q_{i-1}\theta$ are dense.  More formally, make a sequence of such
finite block choices for which all new ratios tend to infinity and
all $\alpha_i\to\lambda_i$.  The Haar law of
$(u_i(x))_{i\in J_j}$ then converges weakly to product Haar measure.
For each fixed nonzero integer character vector, the last nonzero
coefficient times its frequency dominates the sum of the earlier
terms, once the ratios are sufficiently large.  That Fourier
coefficient is therefore exactly zero eventually.  Trigonometric
approximation proves the weak convergence.  The joint law including
the deterministic phases converges to product Haar in the coordinates
and the fixed vector of trial phases.

For fixed $m$ the finitely many squared sampling tests are bounded
functions of these coordinates and phases.  Their discontinuities
have zero limiting measure: every selected $\{n\lambda_i\}$ avoids
the phase cut, and for that phase only finitely many coordinate values
can lie on a cutoff or interval boundary.  Weak convergence therefore
gives convergence of their integrals.  This familiar extension from
continuous tests follows, for example, by continuous upper and lower
approximations outside open neighborhoods of the null discontinuity
sets.  It applies to the squared differences from one as well.
Choose a sufficiently late finite block choice.  Its sampling norms
are below $1/j$, all the strict spread and angle bounds persist, and
all frequency ratios determined so far satisfy their required bounds.
The ratio connecting
the previous block to this one can be made large enough to meet that
block's outstanding condition; the ratio at the end of this block
is imposed at the next stage.  This completes the induction.
\end{proof}

Two consequences of the construction will be used repeatedly:
\begin{equation}\label{eq:tail-sizes}
 \sum_{s>j}\frac1{K_s}\le\frac2{K_{j+1}},\qquad
 \sum_j\frac{1+h_j^*}{K_j}<\infty,\qquad
 \frac{Q_j}{K_{j+1}}\longrightarrow0.
\end{equation}
They follow respectively from $K_{j+1}>2K_j$, the final inequality in
\eqref{eq:scales}, and $7/10<3/4$.
The parameters are now fixed independently of any putative invariant
projection.  We next turn their sampling identities into operator
identities between pairs of points of the circle.
